\chapter{Conclusion}
\label{ch:conclusion}

This project set out to answer a question the paper itself leaves open:
does eALS's parallelizability, argued on paper, hold up when actually
built in a distributed engine? The answer, based on the experiments in
Chapters~\ref{ch:results} and~\ref{ch:scalability}, is yes. The PySpark
implementation reproduces the paper's element-wise update rule term for
term (Chapter~\ref{ch:methods}), reaches the same optimal missing-data
weight $c_0=512$ the paper reports for the Amazon dataset despite training
on a differently-sized snapshot of it, converges to a stable solution
within roughly ten iterations with no learning rate to tune, and scales in
training time with $10.9\times$ growth for a $16\times$ increase in $K$,
far below the $4096\times$ a cubic solver would need, and it lands in the
same qualitative regime as Spark MLlib's own production-grade ALS.

\medskip
Not everything reproduced cleanly. The popularity-aware weighting
scheme's advantage over plain uniform weighting, a clearly significant
effect in the paper, showed up here only as a small and statistically
untested difference (Section~\ref{sec:weighting}), most plausibly because
this project's dataset, filtered from a different and smaller snapshot of
the same source, carries a less pronounced popularity signal than the
one the paper trained on. The parallelism experiment's finding that four
partitions beat eight or sixteen (Section~\ref{sec:scale_par}) is a
statement about single-machine task-scheduling overhead, not about the
algorithm's parallelizability itself, which Section~\ref{sec:parallel}'s
independence argument guarantees regardless of partition count. And the
project deliberately stops at the paper's offline Algorithm~1, leaving its
online incremental-update extension (Algorithm~2) out of scope, in line
with what the assignment brief asks for.

\medskip
Building the implementation also surfaced a lesson that the paper's own
complexity analysis does not make visible. At this project's scale, the
gap between eALS and Spark MLlib's ALS was not decided by which algorithm
does asymptotically less work, but by an engineering choice: broadcasting
$Q$ and $S^q$ once per sweep instead of shuffling factor updates through
Spark's DataFrame machinery, and a plain Python loop instead of a
vectorized or compiled one (Section~\ref{sec:weaknesses}). A correct
distributed algorithm and a fast distributed implementation of it are two
different things to get right, and this project only fully closed the gap
on the first.

\medskip
The assignment brief really asks two questions: is eALS efficient, and is
it parallelizable without approximation? Both held up under direct
measurement, on real data, in a real distributed engine, which is the one
thing the paper's own single-threaded benchmarks could not show.
